% Laget automatisk av elevark/lag_elevark.py fra sim/fartsvelger.html. Endringer her blir overskrevet.
\documentclass[11pt]{article}
\usepackage{elevark}
\begin{document}
\elevarkhode{Fysikk 2 · Magnetiske felt og krefter}{Fartsvelger}{https://lektorbo.no/sim/fartsvelger.html}{Simuleringen}

Mellom to plater virker det både et elektrisk og et magnetisk felt. Den elektriske kraften og den magnetiske kraften peker hver sin vei. Bare partikler med akkurat riktig fart går rett fram og kommer gjennom spalten på den andre siden.

\section*{Oppgaver}

Start med et proton, akselerasjonsspenning 800~V, $B = 50\,\mathrm{mT}$ og 300~V mellom platene. $e = 1{,}60 \cdot 10^{-19}\,\mathrm{C}$ og $m_{\mathrm{p}} = 1{,}67 \cdot 10^{-27}\,\mathrm{kg}$.

\begin{oppgaveliste}
\oppgave \fase{gjett}{Gjett}Tegn kreftene på protonet når det kommer inn mellom platene. Hvilken kraft er størst? Hvilken plate tror du protonet treffer? Test.
\svar
\oppgave \fase{test}{Regn}Regn ut farten til protonet. Hvor stor må spenningen mellom platene være for at protonet skal gå rett fram? Skriv inn svaret og test.
\svar
\oppgave \fase{forklar}{Forklar}Vis at partikler med farten $v = E/B$ går rett fram, uansett masse og ladning. Bytt til et elektron, og juster akselerasjonsspenningen slik at det får samme fart som protonet. Går elektronet rett fram? Hvorfor, når begge kreftene har snudd?
\svar
\oppgave \fase{gjett}{Gjett}Slå på «Send en stråle med ulike farter». Hvilke partikler bøyes opp, og hvilke bøyes ned? Forklar ut fra hvordan $F_{\mathrm{E}}$ og $F_{\mathrm{B}}$ avhenger av farten.
\svar
\oppgave \fase{test}{Undersøk}Bytt til en alfapartikkel med de samme innstillingene. Hvorfor kommer den ikke gjennom? Hvilken spenning må du ha mellom platene for at den skal komme gjennom?
\svar
\oppgave \fase{test}{Regn}Et elektron akselereres gjennom 200~V og sendes inn i fartsvelgeren med $B = 0{,}50\,\mathrm{mT}$. Hvor stor må det elektriske feltet og spenningen mellom platene være? Test. Hvorfor blir avbøyningen så liten når spenningen er litt feil?
\svar
\oppgave \fase{utfordring}{Utfordring}Etter fartsvelgeren i oppgave 2 går protonene inn i et nytt magnetfelt på 0,10~T, slik som i et massespektrometer. Hvor stor blir radien i banen? Hvorfor er det nyttig å ha en fartsvelger foran et massespektrometer?
\svar
\end{oppgaveliste}

\end{document}
